\documentclass[12pt,a4paper,openany,oneside]{book}

% --- Paquetes ---
\usepackage[utf8]{inputenc}
\usepackage[spanish, es-noshorthands]{babel}
\usepackage{amsmath, amsfonts, amssymb}
\usepackage{graphicx}
\usepackage{geometry}
\usepackage{hyperref}
\usepackage{booktabs}
\usepackage{fancyhdr}
\usepackage{listings}
\usepackage{xcolor}
\usepackage{tikz}
\usepackage{pgfplots}
\usepackage{colortbl}
\usepackage{multirow}
\usepackage{hhline}
\usetikzlibrary{arrows.meta, positioning, shapes.geometric, calc, shapes, chains, scopes}
\pgfplotsset{compat=1.18}

\geometry{margin=2.5cm}

% --- Colores y estilos para código Python ---
\definecolor{codegreen}{rgb}{0,0.6,0}
\definecolor{codegray}{rgb}{0.5,0.5,0.5}
\definecolor{codepurple}{rgb}{0.58,0,0.82}
\definecolor{backcolour}{rgb}{0.95,0.95,0.92}

\lstset{
    language=Python,
    backgroundcolor=\color{backcolour},   
    commentstyle=\color{codegreen},
    keywordstyle=\color{blue},
    numberstyle=\tiny\color{codegray},
    stringstyle=\color{codepurple},
    basicstyle=\ttfamily\small,
    breaklines=true,
    frame=single,
    showstringspaces=false
}

% --- Estilo de cabecera y pie ---
\pagestyle{fancy}
\fancyhf{}
\renewcommand{\headrulewidth}{0.4pt}
\renewcommand{\footrulewidth}{0.4pt}
\fancyhead[L]{\small Carlos César Sánchez Coronel}
\fancyhead[R]{\small Sesión 7: Redes Recurrentes para NLP}
\fancyfoot[L]{\small Carlos César Sánchez Coronel}
\fancyfoot[C]{\thepage}

% --- Portada ---
\title{
    \textbf{Sesión 7} \\ 
    \Huge \textbf{REDES RECURRENTES (RNN, LSTM, GRU) PARA NLP} \\
    \Large Modelado de secuencias y memoria a largo plazo
}
\author{
    \small Por \\ \vspace{0.2cm}
    \Large \textsc{Carlos César Sánchez Coronel}
}
\date{2026}

\begin{document}

\maketitle
\tableofcontents
\addcontentsline{toc}{chapter}{Índice general}

% ------------------------------------------------------------------
% CAPÍTULO 7: SESIÓN 7 - REDES RECURRENTES PARA NLP
% ------------------------------------------------------------------
\chapter{Redes Recurrentes (RNN, LSTM, GRU) para NLP}

En la sesión anterior se construyeron redes feed-forward que procesaban oraciones como un todo (promediando embeddings). Sin embargo, el lenguaje es inherentemente secuencial: el significado de una palabra depende de las palabras anteriores. Las redes recurrentes (RNN) están diseñadas para modelar secuencias, manteniendo un estado oculto que se actualiza con cada nuevo elemento. Esta sesión cubre las RNN básicas, sus limitaciones (gradiente desvaneciente), y las arquitecturas que las superan: LSTM y GRU. Se incluyen fundamentos matemáticos, implementaciones en PyTorch y aplicaciones prácticas.

\section{Logro de la sesión}
Modelar secuencias de texto utilizando redes recurrentes, comprender el problema del gradiente desvaneciente y conocer las variantes LSTM y GRU que lo mitigan, así como sus aplicaciones en tareas de NLP.

\section{Introducción: la necesidad de modelar secuencias}
Las redes feed-forward procesan cada entrada de forma independiente. Para texto, esto significa que se pierde el orden de las palabras. Consideremos dos oraciones:
\begin{itemize}
    \item "El gato persigue al ratón"
    \item "El ratón persigue al gato"
\end{itemize}
Una red feed-forward que promedia embeddings produciría el mismo vector para ambas, a pesar de tener significados opuestos. Las redes recurrentes resuelven esto procesando la secuencia palabra por palabra, manteniendo un estado que acumula información del pasado.

\section{Redes Neuronales Recurrentes (RNN)}

\subsection{Arquitectura básica}
Una RNN procesa una secuencia de vectores \(\mathbf{x}_1, \mathbf{x}_2, ..., \mathbf{x}_T\) (por ejemplo, embeddings de palabras). En cada paso \(t\), la red actualiza un estado oculto \(\mathbf{h}_t\) basado en el estado anterior \(\mathbf{h}_{t-1}\) y la entrada actual \(\mathbf{x}_t\):

\begin{equation}
\mathbf{h}_t = \tanh( \mathbf{W}_{hh} \mathbf{h}_{t-1} + \mathbf{W}_{xh} \mathbf{x}_t + \mathbf{b}_h )
\end{equation}

Luego, puede producir una salida \(\mathbf{y}_t\) (por ejemplo, una predicción) en cada paso:

\begin{equation}
\mathbf{y}_t = \text{softmax}( \mathbf{W}_{hy} \mathbf{h}_t + \mathbf{b}_y )
\end{equation}

\begin{center}
\begin{tikzpicture}[
    cell/.style={rectangle, draw, minimum width=1.5cm, minimum height=1cm},
    arrow/.style={-latex, thick}
]
\node[cell] (h0) at (0,0) {$\mathbf{h}_0$};
\node[cell] (h1) at (3,0) {$\mathbf{h}_1$};
\node[cell] (h2) at (6,0) {$\mathbf{h}_2$};
\node[cell] (h3) at (9,0) {$\mathbf{h}_3$};

\node at (0, -1.5) {$\mathbf{x}_1$};
\node at (3, -1.5) {$\mathbf{x}_2$};
\node at (6, -1.5) {$\mathbf{x}_3$};

\draw[arrow] (0, -1) -- (h1);
\draw[arrow] (3, -1) -- (h2);
\draw[arrow] (6, -1) -- (h3);

\draw[arrow] (h0) -- (h1);
\draw[arrow] (h1) -- (h2);
\draw[arrow] (h2) -- (h3);

\node at (1.5, 1) {...
\draw[arrow] (h3.east) -- ++(1,0);
\end{tikzpicture}
\caption{Arquitectura de una RNN desplegada en el tiempo.}
\end{center}

\subsection{Entrenamiento: Backpropagation Through Time (BPTT)}
El gradiente se propaga a través de los pasos de tiempo. La pérdida total es la suma de las pérdidas en cada paso (o solo en el último, según la tarea). El gradiente respecto a los pesos involucra productos de matrices, lo que puede causar problemas de estabilidad.

\subsection{Problema del gradiente desvaneciente/explosivo}
Durante la retropropagación, el gradiente se multiplica repetidamente por la matriz de pesos recurrente \(\mathbf{W}_{hh}\). Si los autovalores de \(\mathbf{W}_{hh}\) son menores que 1, el gradiente tiende a cero (desvaneciente); si son mayores, el gradiente explota. Esto impide que la RNN aprenda dependencias a largo plazo.

\textbf{Ejemplo}: En una oración larga, la palabra "gato" al principio puede influir en un verbo al final, pero el gradiente desvaneciente hace que esa influencia sea insignificante.

\section{LSTM (Long Short-Term Memory)}

LSTM fue propuesta por Hochreiter \& Schmidhuber (1997) para resolver el problema del gradiente desvaneciente mediante una arquitectura con celdas de memoria y puertas que controlan el flujo de información.

\subsection{Arquitectura de una celda LSTM}
Cada celda LSTM en el paso \(t\) recibe:
\begin{itemize}
    \item Entrada \(\mathbf{x}_t\) (embedding de la palabra actual)
    \item Estado oculto anterior \(\mathbf{h}_{t-1}\)
    \item Estado de celda anterior \(\mathbf{c}_{t-1}\) (memoria a largo plazo)
\end{itemize}
Produce:
\begin{itemize}
    \item Estado oculto actual \(\mathbf{h}_t\)
    \item Estado de celda actual \(\mathbf{c}_t\)
\end{itemize}

\begin{center}
\begin{tikzpicture}[
    cell/.style={rectangle, draw, minimum width=2.5cm, minimum height=2cm},
    arrow/.style={-latex, thick}
]
\node[cell] (lstm) at (0,0) {LSTM};
\node[left=1cm of lstm] (x) {$\mathbf{x}_t$};
\node[above left=0.5cm of lstm] (hprev) {$\mathbf{h}_{t-1}$};
\node[below left=0.5cm of lstm] (cprev) {$\mathbf{c}_{t-1}$};
\node[above right=0.5cm of lstm] (h) {$\mathbf{h}_t$};
\node[below right=0.5cm of lstm] (c) {$\mathbf{c}_t$};
\draw[arrow] (x) -- (lstm);
\draw[arrow] (hprev) -- (lstm);
\draw[arrow] (cprev) -- (lstm);
\draw[arrow] (lstm) -- (h);
\draw[arrow] (lstm) -- (c);
\end{tikzpicture}
\caption{Entradas y salidas de una celda LSTM.}
\end{center}

\subsection{Puertas y ecuaciones}

\subsubsection{Puerta de olvido (forget gate)}
Decide qué información del estado de celda anterior debe descartarse:
\begin{equation}
\mathbf{f}_t = \sigma( \mathbf{W}_f [\mathbf{h}_{t-1}, \mathbf{x}_t] + \mathbf{b}_f )
\end{equation}
donde \(\sigma\) es la función sigmoide, que produce valores entre 0 y 1.

\subsubsection{Puerta de entrada (input gate)}
Decide qué nueva información se almacenará en el estado de celda:
\begin{equation}
\mathbf{i}_t = \sigma( \mathbf{W}_i [\mathbf{h}_{t-1}, \mathbf{x}_t] + \mathbf{b}_i )
\end{equation}

\subsubsection{Candidato a nueva memoria}
\begin{equation}
\tilde{\mathbf{c}}_t = \tanh( \mathbf{W}_c [\mathbf{h}_{t-1}, \mathbf{x}_t] + \mathbf{b}_c )
\end{equation}

\subsubsection{Actualización del estado de celda}
Se combina el olvido del estado anterior y la nueva memoria filtrada por la puerta de entrada:
\begin{equation}
\mathbf{c}_t = \mathbf{f}_t \odot \mathbf{c}_{t-1} + \mathbf{i}_t \odot \tilde{\mathbf{c}}_t
\end{equation}
donde \(\odot\) denota producto elemento a elemento (Hadamard).

\subsubsection{Puerta de salida (output gate)}
Decide qué parte del estado de celda se emite como estado oculto:
\begin{equation}
\mathbf{o}_t = \sigma( \mathbf{W}_o [\mathbf{h}_{t-1}, \mathbf{x}_t] + \mathbf{b}_o )
\end{equation}

\subsubsection{Estado oculto final}
\begin{equation}
\mathbf{h}_t = \mathbf{o}_t \odot \tanh(\mathbf{c}_t)
\end{equation}

\subsection{Interpretación}
\begin{itemize}
    \item La puerta de olvido permite que la celda olvide información irrelevante.
    \item La puerta de entrada decide qué nueva información es relevante.
    \item El estado de celda \(\mathbf{c}_t\) actúa como una autopista de gradiente: las operaciones son lineales (sumas y productos elemento a elemento), lo que facilita que el gradiente fluya sin desvanecerse.
    \item La puerta de salida controla cuánto de la memoria se expone al resto de la red.
\end{itemize}

\subsection{Ventajas de LSTM}
- Maneja dependencias a largo plazo (cientos de pasos).
- El gradiente puede fluir a través de \(\mathbf{c}_t\) con pocas transformaciones no lineales.
- Es el estado del arte en muchas tareas secuenciales antes de los Transformers.

\section{GRU (Gated Recurrent Unit)}

GRU, propuesta por Cho et al. (2014), es una simplificación de LSTM que combina las puertas de olvido y entrada en una sola "puerta de actualización", y fusiona el estado de celda y el estado oculto.

\subsection{Arquitectura GRU}
En cada paso \(t\):
\begin{itemize}
    \item Entrada \(\mathbf{x}_t\)
    \item Estado oculto anterior \(\mathbf{h}_{t-1}\)
    \item Produce nuevo estado oculto \(\mathbf{h}_t\)
\end{itemize}

\subsubsection{Puerta de reinicio (reset gate)}
Decide cuánto del estado anterior se olvida:
\begin{equation}
\mathbf{r}_t = \sigma( \mathbf{W}_r [\mathbf{h}_{t-1}, \mathbf{x}_t] + \mathbf{b}_r )
\end{equation}

\subsubsection{Puerta de actualización (update gate)}
Decide cuánto del nuevo estado se copia del anterior:
\begin{equation}
\mathbf{z}_t = \sigma( \mathbf{W}_z [\mathbf{h}_{t-1}, \mathbf{x}_t] + \mathbf{b}_z )
\end{equation}

\subsubsection{Candidato a nuevo estado}
\begin{equation}
\tilde{\mathbf{h}}_t = \tanh( \mathbf{W}_h [\mathbf{r}_t \odot \mathbf{h}_{t-1}, \mathbf{x}_t] + \mathbf{b}_h )
\end{equation}

\subsubsection{Actualización del estado oculto}
\begin{equation}
\mathbf{h}_t = (1 - \mathbf{z}_t) \odot \mathbf{h}_{t-1} + \mathbf{z}_t \odot \tilde{\mathbf{h}}_t
\end{equation}

\subsection{Ventajas de GRU}
- Tiene menos parámetros que LSTM (tres puertas en lugar de cuatro), por lo que es más rápido de entrenar y menos propenso a sobreajuste con datos pequeños.
- En muchos problemas, su rendimiento es comparable al de LSTM.
- Es más fácil de implementar y depurar.

\subsection{Comparación LSTM vs GRU}
\begin{table}[h]
\centering
\begin{tabular}{|l|c|c|}
\hline
\textbf{Característica} & \textbf{LSTM} & \textbf{GRU} \\
\hline
Puertas & 3 (input, forget, output) + celda & 2 (update, reset) \\
Estado de celda separado & Sí & No (solo estado oculto) \\
Parámetros & Más & Menos \\
Rendimiento en datos grandes & Excelente & Muy bueno \\
Rendimiento en datos pequeños & Puede sobreajustar & Más robusto \\
\hline
\end{tabular}
\caption{Comparación entre LSTM y GRU.}
\end{table}

\section{Aplicaciones de RNN/LSTM/GRU en NLP}

\subsection{Clasificación de sentimiento}
Se toma el último estado oculto (o un promedio de todos) y se pasa por una capa fully-connected para clasificar. Ejemplo: determinar si una reseña es positiva o negativa.

\subsection{Generación de texto}
Se entrena una RNN para predecir la siguiente palabra dada la secuencia anterior. Luego se usa para generar texto nuevo muestreando de la distribución de salida. Ejemplo: generación de poesía, código, etc.

\subsection{Traducción automática}
Modelos seq2seq (encoder-decoder) donde un encoder RNN procesa la frase origen y un decoder RNN genera la frase destino. La atención (próxima sesión) mejora este proceso.

\subsection{Reconocimiento de entidades nombradas (NER)}
Se etiqueta cada palabra con su categoría (persona, lugar, etc.) usando la salida de cada paso.

\subsection{Análisis de sentimiento a nivel de aspecto}
Determinar la polaridad hacia aspectos específicos (ej. "la batería dura mucho" → aspecto "batería", positivo).

\section{Ejemplo paso a paso: Clasificación de sentimiento con LSTM en PyTorch}

\subsection{Datos}
Usaremos el mismo conjunto de reseñas pequeñas de la sesión anterior, pero ahora con secuencias más largas para ilustrar la capacidad de LSTM.

\begin{lstlisting}[language=Python]
import torch
import torch.nn as nn
import torch.optim as optim
from collections import Counter
import numpy as np

# Datos de ejemplo (ampliados)
texts = [
    "I loved this movie, it was absolutely fantastic and I will watch it again",
    "Terrible film, I hated it, worst movie ever, waste of time",
    "Amazing acting and great story, the characters are very well developed",
    "Boring and slow, nothing happens, completely waste of money"
]
labels = [1, 0, 1, 0]
\end{lstlisting}

\subsection{Paso 1: Construir vocabulario y codificar}
\begin{lstlisting}[language=Python]
def build_vocab(texts, max_size=10000):
    word_counts = Counter()
    for text in texts:
        word_counts.update(text.lower().split())
    vocab = {word: idx+2 for idx, (word, _) in enumerate(word_counts.most_common(max_size))}
    vocab['<PAD>'] = 0
    vocab['<UNK>'] = 1
    return vocab

vocab = build_vocab(texts)
vocab_size = len(vocab)

def encode(text, vocab, max_len=20):
    tokens = text.lower().split()
    indices = [vocab.get(token, vocab['<UNK>']) for token in tokens[:max_len]]
    # padding
    indices += [vocab['<PAD>']] * (max_len - len(indices))
    return indices

X = [encode(text, vocab) for text in texts]
X = torch.tensor(X)  # (4, 20)
y = torch.tensor(labels)
\end{lstlisting}

\subsection{Paso 2: Definir modelo LSTM}
\begin{lstlisting}[language=Python]
class LSTMSentiment(nn.Module):
    def __init__(self, vocab_size, embed_dim=50, hidden_dim=64, num_layers=1, num_classes=2, dropout=0.5):
        super().__init__()
        self.embedding = nn.Embedding(vocab_size, embed_dim, padding_idx=0)
        self.lstm = nn.LSTM(embed_dim, hidden_dim, num_layers, batch_first=True, dropout=dropout)
        self.fc = nn.Linear(hidden_dim, num_classes)
        self.dropout = nn.Dropout(dropout)
        
    def forward(self, x):
        # x: (batch, seq_len)
        embedded = self.embedding(x)  # (batch, seq_len, embed_dim)
        lstm_out, (hidden, cell) = self.lstm(embedded)
        # Usamos el último estado oculto (de la última capa, último paso)
        last_hidden = hidden[-1]  # (batch, hidden_dim)
        out = self.fc(self.dropout(last_hidden))
        return out

model = LSTMSentiment(vocab_size, embed_dim=50, hidden_dim=64, num_layers=1)
\end{lstlisting}

\subsection{Paso 3: Entrenamiento}
\begin{lstlisting}[language=Python]
criterion = nn.CrossEntropyLoss()
optimizer = optim.Adam(model.parameters(), lr=0.001, weight_decay=1e-5)

epochs = 200
for epoch in range(epochs):
    optimizer.zero_grad()
    logits = model(X)
    loss = criterion(logits, y)
    loss.backward()
    optimizer.step()
    
    if (epoch+1) % 50 == 0:
        print(f'Epoch {epoch+1}, Loss: {loss.item():.4f}')
\end{lstlisting}

\subsection{Paso 4: Evaluación}
\begin{lstlisting}[language=Python]
with torch.no_grad():
    logits = model(X)
    predictions = torch.argmax(logits, dim=1)
    accuracy = (predictions == y).float().mean()
    print(f'Accuracy: {accuracy:.2%}')
\end{lstlisting}

\section{Métricas en tareas de NLP con RNN}

\subsection{Perplejidad (Perplexity)}
Especialmente para modelos de lenguaje (generación), la perplejidad mide qué tan "sorprendido" está el modelo por los datos. Se define como:
\[
\text{PPL} = \exp\left( -\frac{1}{N} \sum_{i=1}^N \log P(w_i | \text{contexto}) \right)
\]
Es el exponencial de la entropía cruzada. Valores más bajos indican mejor modelo.

\subsection{Accuracy}
Para clasificación de secuencias completas.

\subsection{F1-score, precisión, recall}
Para tareas de etiquetado a nivel de palabra (NER, POS tagging), se calculan sobre cada token (con promedios macro/micro).

\subsection{Exactitud de secuencia (Sequence Accuracy)}
En tareas de generación, se puede medir si la secuencia generada coincide exactamente con la objetivo (raro en práctica).

\section{Limitaciones de RNN/LSTM}

\subsection{Computacionalmente secuenciales}
No pueden paralelizarse a lo largo de la dimensión temporal, lo que los hace lentos en GPUs en comparación con Transformers.

\subsection{Dificultad con dependencias muy largas}
Aunque LSTM mejora, aún puede tener problemas con secuencias de miles de pasos.

\subsection{Memoria limitada}
El estado oculto tiene capacidad finita; no puede almacenar toda la información de una secuencia larga.

\section{Preparación para entrevistas}

\subsection{Preguntas típicas}
\begin{itemize}
    \item \textbf{"Explica por qué LSTM funciona mejor que RNN vanilla para secuencias largas."}
    Respuesta: RNN sufre de gradiente desvaneciente debido a la multiplicación repetida de la matriz de pesos. LSTM introduce una celda de memoria con puertas que permiten que el gradiente fluya a través de operaciones lineales (sumas y productos elemento a elemento), preservándolo a lo largo del tiempo.
    
    \item \textbf{"¿Cuál es la diferencia entre LSTM y GRU?"}
    GRU tiene dos puertas (update y reset) y no tiene estado de celda separado; combina la memoria en el estado oculto. Tiene menos parámetros y puede ser más rápido, pero LSTM puede capturar dependencias más complejas con su celda de memoria.
    
    \item \textbf{"¿Qué problema resuelve la puerta de olvido en LSTM?"}
    Permite que la red decida qué información del pasado es irrelevante y debe descartarse, evitando que la celda se sature con información inútil.
    
    \item \textbf{"¿Cómo se entrena una RNN?"}
    Con backpropagation through time (BPTT), que propaga el gradiente a través de los pasos de tiempo. Se usa truncamiento para secuencias muy largas.
    
    \item \textbf{"¿Cuándo usarías una RNN en lugar de una red feed-forward?"}
    Cuando los datos son secuenciales y el orden importa: texto, series temporales, audio, video. Si el orden no importa, una feed-forward puede ser suficiente.
\end{itemize}

\subsection{Preguntas avanzadas}
\begin{itemize}
    \item \textbf{"¿Qué es el truncamiento en BPTT?"}
    Se limita la propagación del gradiente a un número fijo de pasos hacia atrás para evitar problemas de memoria y cómputo, especialmente en secuencias muy largas.
    
    \item \textbf{"¿Cómo manejarías secuencias de longitud variable en un lote?"}
    Usando padding (rellenar con un token especial) y máscaras para que la RNN ignore esos pasos. En PyTorch, se usa `pack_padded_sequence` para eficiencia.
    
    \item \textbf{"¿Qué son las RNN bidireccionales?"}
    Procesan la secuencia en ambas direcciones (izquierda a derecha y derecha a izquierda), concatenando los estados ocultos. Útiles cuando se tiene acceso a toda la secuencia (ej. clasificación de oraciones completas).
\end{itemize}

\section{Conexión con el resto del curso}
\begin{itemize}
    \item \textbf{Sesión 6 (Feed-Forward)}: Las RNN reemplazan el promedio de embeddings con un procesamiento secuencial.
    \item \textbf{Sesión 8 (Atención y Transformers)}: La atención resuelve las limitaciones de las RNN al permitir acceso directo a cualquier paso.
    \item \textbf{Sesiones anteriores}: Las capas de embedding y regularización son comunes.
\end{itemize}

\section{Resumen}
\begin{itemize}
    \item Las \textbf{RNN} procesan secuencias manteniendo un estado oculto, pero sufren de gradiente desvaneciente.
    \item \textbf{LSTM} introduce puertas (olvido, entrada, salida) y un estado de celda para preservar gradientes y capturar dependencias a largo plazo.
    \item \textbf{GRU} simplifica LSTM combinando puertas y reduciendo parámetros.
    \item Se aplican a clasificación de texto, generación, traducción, NER, etc.
    \item Aunque han sido superados por Transformers, siguen siendo útiles en muchos escenarios y son fundamentales en la historia del NLP.
\end{itemize}

\end{document}